% ============================================================
%  meta.tex —— 课程元信息（整份试卷唯一需要改动的配置）
%  换一门课，只改这个文件，style.tex 不用动
% ============================================================

% ---- 学校与课程 ----
\newcommand{\SchoolName}{××大学}
\newcommand{\CollegeName}{}          % 学院名，留空则不打印；填了会接在学校名后面
\newcommand{\CourseName}{信号与系统}
\newcommand{\TermName}{2026—2027 学年\quad 第 1 学期}
\newcommand{\ExamDate}{}             % 考试日期，留空则不打印

% ---- 卷面抬头 ----
\newcommand{\MajorList}{电子信息工程、通信工程、自动化}   % 适用专业
\newcommand{\ExamForm}{闭卷}                              % 开卷 / 闭卷 / 半开卷
\newcommand{\ExamTime}{120 分钟}
\newcommand{\TotalPoints}{100}
\newcommand{\ScopeNote}{全课程}                           % 命题范围

% ---- 卷面要素开关 ----
% 1 = 显示，0 = 隐藏。密封线在双面印刷或不要求装订时可关掉
\newcommand{\ShowSealLine}{1}
\newcommand{\ShowScoreTable}{1}

% ---- 密封线文字（仅 \ShowSealLine=1 时生效）----
\newcommand{\SealText}{%
班级\underline{\hspace{1.5cm}}\quad 学号\underline{\hspace{1.75cm}}\quad 姓名\underline{\hspace{1.45cm}}\qquad
{\heiti 密}\ {\heiti 封}\ {\heiti 线}\qquad 请勿在密封线内答题}

% ---- 试卷说明（三到四条为宜，超出会挤掉正文空间）----
\newcommand{\ExamNotice}{%
1. 答案一律写在试卷相应位置，写在草稿纸上无效；%
2. 计算题须写出主要公式与演算过程，只有结果不给分；%
3. 答卷前请将密封线内项目填写清楚。}

% ---- 答案册评分说明 ----
\newcommand{\GradeNotice}{%
填空题按空给分，答案正确即得分；选择题、判断题不倒扣分；%
简答题按要点给分，答出关键术语与结论即可，表述完整、条理清楚者可在该题分值内适当调整；%
计算题按步骤给分，公式正确而结果错误的扣 1～2 分，只有结果无过程的不得分。}
